\chapter{Complexity — Words}
%%% generated by the blueprint pipeline from TCSlib.Complexity.ClassP.BinaryArith.Words
\section{Overview}

Declarations of the Lean module \texttt{TCSlib.Complexity.ClassP.BinaryArith.Words}.

\section{Declarations}

\begin{definition}[Encoding a list of words]
\lean{Complexity.BinArith.encWords}
\label{Complexity.BinArith.encWords}
\leanok
\uses{Turing.pairEncode}
For a list $L$ of Boolean words, its encoding is the concatenation, in order, of the
pair encodings of each word in $L$ with the empty word.
\end{definition}

\begin{theorem}[Appending after an empty second component]
\lean{Complexity.BinArith.pairEncode_nil_append}
\label{Complexity.BinArith.pairEncode_nil_append}
\leanok
\uses{Turing.pairEncode}
\difficulty{1}
For bit strings $a$ and $b$, appending $b$ to the encoding of the pair $(a,\varnothing)$
gives the encoding of $(a,b)$, where $\varnothing$ denotes the empty bit string.
\end{theorem}

\begin{theorem}[Encoding the empty list]
\lean{Complexity.BinArith.encWords_nil}
\label{Complexity.BinArith.encWords_nil}
\leanok
\uses{Complexity.BinArith.encWords}
\difficulty{1}
The encoding of the empty list of bit strings is the empty bit string.
\end{theorem}

\begin{theorem}[Encoding a list of bit strings]
\lean{Complexity.BinArith.encWords_cons}
\label{Complexity.BinArith.encWords_cons}
\leanok
\uses{Complexity.BinArith.encWords, Complexity.BinArith.pairEncode_nil_append, Turing.pairEncode}
\difficulty{1}
For any bit string $w$ and any finite list $L$ of bit strings, encoding the list
beginning with $w$ gives the pairing of $w$ with the encoding of the remaining list.
\end{theorem}

\begin{theorem}[Encoding a list with one appended word]
\lean{Complexity.BinArith.encWords_append_singleton}
\label{Complexity.BinArith.encWords_append_singleton}
\leanok
\uses{Complexity.BinArith.encWords, Turing.pairEncode}
\difficulty{1}
For every list $L$ of bit strings and every bit string $w$, encoding the list obtained
by appending $w$ to $L$ gives the encoding of $L$ followed by the self-delimiting pair
encoding of $w$ and the empty string.
\end{theorem}

\begin{theorem}[Length of an encoded list of words]
\lean{Complexity.BinArith.length_encWords}
\label{Complexity.BinArith.length_encWords}
\leanok
\uses{Complexity.BinArith.encWords, Turing.length_pairEncode}
\difficulty{4}
For any finite list $L$ of finite bit strings, concatenating the self-delimiting
encodings of each word in $L$ paired with the empty word produces a bit string of length
exactly $\sum_{w\in L}(2|w|+2)$.
\end{theorem}

\begin{theorem}[Bound on encoded list length]
\lean{Complexity.BinArith.length_encWords_le}
\label{Complexity.BinArith.length_encWords_le}
\leanok
\uses{Complexity.BinArith.encWords, Complexity.BinArith.encWords_cons, Turing.length_pairEncode}
\difficulty{4}
Let $L$ be a finite list of finite bit strings, and let $c$ be a natural number such
that every word $w$ in $L$ has length at most $c$. The concatenation of the
self-delimiting encodings of each word paired with the empty word has length at most
$|L|(2c+2)$.
\end{theorem}

\begin{theorem}[Strict length decrease under second projection]
\lean{Complexity.BinArith.length_pairSndD_lt}
\label{Complexity.BinArith.length_pairSndD_lt}
\leanok
\uses{Complexity.pairSndD, Complexity.pairSndD_pairEncode, Turing.eq_pairEncode_of_pairDecode, Turing.length_pairEncode, Turing.pairDecode}
\difficulty{4}
For every nonempty bit string $w$, the total second projection of $w$ has length
strictly less than the length of $w$.
\end{theorem}

\begin{definition}[Decoding a word into words]
\lean{Complexity.BinArith.decWords}
\label{Complexity.BinArith.decWords}
\leanok
\uses{Complexity.BinArith.length_pairSndD_lt, Complexity.pairFstD, Complexity.pairSndD}
For a Boolean word $w$, repeatedly split the current word into its first and second
components, append the first component to the output, and continue with the second
component; stop when the current word is empty. The output is the resulting list of
Boolean words.
\end{definition}

\begin{theorem}[Decoding the empty word]
\lean{Complexity.BinArith.decWords_nil}
\label{Complexity.BinArith.decWords_nil}
\leanok
\uses{Complexity.BinArith.decWords}
\difficulty{3}
The recursive decoder of a bit string into a list of words sends the empty word to the
empty list.
\end{theorem}

\begin{theorem}[Decoder recursion on nonempty words]
\lean{Complexity.BinArith.decWords_of_ne}
\label{Complexity.BinArith.decWords_of_ne}
\leanok
\uses{Complexity.BinArith.decWords, Complexity.pairFstD, Complexity.pairSndD}
\difficulty{3}
For every nonempty bit string $w$, its decoded list of words is the first projection of
the pair parser applied to $w$, followed by the decoded list of words obtained from the
second projection of that parse.
\end{theorem}

\begin{theorem}[Decoding an encoded list of words]
\lean{Complexity.BinArith.decWords_encWords}
\label{Complexity.BinArith.decWords_encWords}
\leanok
\uses{Complexity.BinArith.decWords, Complexity.BinArith.decWords_of_ne, Complexity.BinArith.encWords, Complexity.BinArith.encWords_cons, Complexity.pairFstD_pairEncode, Complexity.pairSndD_pairEncode, Turing.pairEncode}
\difficulty{4}
For every finite list $L$ of finite bit strings, encode each word $x$ as the
self-delimiting pair of $x$ with the empty bit string, and concatenate these encodings.
Applying the recursive word-list decoder to the resulting bit string recovers $L$.
\end{theorem}

\begin{theorem}[Bound on recursive pair decomposition]
\lean{Complexity.BinArith.length_decWords_le}
\label{Complexity.BinArith.length_decWords_le}
\leanok
\uses{Complexity.BinArith.decWords, Complexity.BinArith.decWords_of_ne, Complexity.BinArith.length_pairSndD_lt, Complexity.pairSndD}
\difficulty{4}
Starting from a bit string $w$, form a list recursively: use the empty list when $w$ is
empty; otherwise prepend the first component obtained by parsing $w$ (using the empty
string if parsing fails) and continue with its parsed second component (again using the
empty string if parsing fails). The resulting list has length at most the length of $w$.
\end{theorem}

\begin{theorem}[Length bound for decoded words]
\lean{Complexity.BinArith.length_le_of_mem_decWords}
\label{Complexity.BinArith.length_le_of_mem_decWords}
\leanok
\uses{Complexity.BinArith.decWords, Complexity.BinArith.decWords_of_ne, Complexity.BinArith.length_pairSndD_lt, Complexity.length_pairFstD_le}
\difficulty{4}
Starting from a bit string $w$, repeatedly output the total first projection of the
current string and continue with its total second projection, stopping when the current
string is empty. Every word $x$ in the resulting list has length at most $w$'s length.
\end{theorem}

\begin{theorem}[Length under iterated second-component decoding]
\lean{Complexity.BinArith.length_iterate_pairSndD_le}
\label{Complexity.BinArith.length_iterate_pairSndD_le}
\leanok
\uses{Complexity.length_pairSndD_le, Complexity.pairSndD}
\difficulty{4}
Let $w$ be any finite bit string, and let $i$ be a natural number. Repeatedly take the
second component returned by the aligned pair parser, using the empty word whenever
parsing fails. After $i$ such operations, the resulting bit string has length at most
$|w|$.
\end{theorem}

\begin{theorem}[Projection iteration on the empty word]
\lean{Complexity.BinArith.iterate_pairSndD_nil}
\label{Complexity.BinArith.iterate_pairSndD_nil}
\leanok
\uses{Complexity.pairSndD}
\difficulty{1}
For every natural number $i$, applying the total second projection $i$ times to the
empty bit string gives the empty bit string.
\end{theorem}

\begin{theorem}[Iterated decoding and concatenation]
\lean{Complexity.BinArith.flatMap_listLoop}
\label{Complexity.BinArith.flatMap_listLoop}
\leanok
\uses{Complexity.BinArith.decWords, Complexity.BinArith.decWords_of_ne, Complexity.BinArith.length_pairSndD_lt, Complexity.pairFstD, Complexity.pairSndD, Turing.pairDecode}
\difficulty{5}
Let $c$ be any function from bit strings to bit strings, let $N$ be a natural number,
and let $w$ be a bit string with $|w|\le N$. For each $i=0,\ldots,N-1$, if the second
decoded component obtained by applying the second projection $i$ times to $w$ is
nonempty, apply $c$ to the first decoded component of that string; concatenate these
outputs. This equals the concatenation of the images under $c$ of all words obtained by
repeatedly decoding $w$.
\end{theorem}

\begin{definition}[One list-processing step]
\lean{Complexity.BinArith.listStep}
\label{Complexity.BinArith.listStep}
\leanok
\uses{Complexity.pairFstD, Complexity.pairSndD, Turing.pairEncode}
For a Boolean word $S$, form the pair encoding of the first component of $S$ and the
second component of the second component of $S$.
\end{definition}

\begin{definition}[Mapping one encoded chunk]
\lean{Complexity.BinArith.mapChunk}
\label{Complexity.BinArith.mapChunk}
\leanok
\uses{Complexity.pairFstD, Complexity.pairSndD, Turing.pairEncode}
Given a function $g$ from Boolean words to Boolean words and a Boolean word $S$, return
the empty word if the second component of $S$ is empty. Otherwise, pair-encode the
result of applying $g$ to the pair encoding of the first component of $S$ and the first
component of its second component, with the empty word.
\end{definition}

\begin{definition}[Mapping across a word]
\lean{Complexity.BinArith.mapW}
\label{Complexity.BinArith.mapW}
\leanok
\uses{Complexity.BinArith.listStep, Complexity.BinArith.mapChunk}
Given a function $g$ from Boolean words to Boolean words and a Boolean word $S$,
concatenate the results of mapping one encoded chunk over the iterates of the
list-processing step applied to $S$, for iteration indices from $0$ through $|S|+1$.
\end{definition}

\begin{theorem}[Iteration of the list step]
\lean{Complexity.BinArith.iterate_listStep}
\label{Complexity.BinArith.iterate_listStep}
\leanok
\uses{Complexity.BinArith.listStep, Complexity.pairSndD, Turing.pairEncode}
\difficulty{4}
For any bit strings $q,w$ and natural number $i$, applying the list step $i$ times to
the pair encoding $\langle q,w\rangle$ gives $\langle q,\pi_2^{[i]}(w)\rangle$, where
$\pi_2$ is the total second projection and $\pi_2^{[i]}$ denotes its $i$-fold iterate.
\end{theorem}

\begin{theorem}[Mapping over encoded words]
\lean{Complexity.BinArith.mapW_pairEncode}
\label{Complexity.BinArith.mapW_pairEncode}
\leanok
\uses{Complexity.BinArith.decWords, Complexity.BinArith.encWords, Complexity.BinArith.flatMap_listLoop, Complexity.BinArith.iterate_listStep, Complexity.BinArith.mapChunk, Complexity.BinArith.mapW, Complexity.pairFstD_pairEncode, Complexity.pairSndD_pairEncode, Turing.length_pairEncode, Turing.pairEncode}
\difficulty{3}
Let $g$ be any function from bit strings to bit strings, and let $q$ and $w$ be bit
strings. Applying the map operation to the encoding of $(q,w)$ yields the concatenated
encodings, with empty second component, of the words obtained by decoding $w$ and
applying $g$ to the encoding of each decoded word.
\end{theorem}

\begin{theorem}[Polynomial-Time Closure Under Word Mapping]
\lean{Complexity.BinArith.polyTimeComputable_mapW}
\label{Complexity.BinArith.polyTimeComputable_mapW}
\leanok
\uses{Complexity.BinArith.iterate_listStep, Complexity.BinArith.length_iterate_pairSndD_le, Complexity.BinArith.listStep, Complexity.BinArith.mapChunk, Complexity.BinArith.mapW, Complexity.BinArith.polyTimeComputable_slot2, Complexity.BinArith.polyTimeComputable_slot3, Complexity.PolyTimeComputable, Complexity.PolyTimeComputable.comp, Complexity.PolyTimeComputable.pairEncode, Complexity.length_pairSndD_le, Complexity.length_pair_components_le, Complexity.pairFstD, Complexity.pairSndD, Complexity.polyTimeComputable_const, Complexity.polyTimeComputable_emitIter, Complexity.polyTimeComputable_isNil, Complexity.polyTimeComputable_ite, Complexity.polyTimeComputable_pairFstD, Complexity.polyTimeComputable_pairSndD, Turing.length_pairEncode, Turing.pairEncode}
\difficulty{5}
If $g$ is a polynomial-time computable function on bit strings, then the following
function is polynomial-time computable: encode the two decoded components of the input
as an initial pair, and repeatedly apply the prescribed pair update for $|z|+2$
iterations, concatenate the resulting mapped chunks, and return the resulting bit
string.
\end{theorem}

\begin{definition}[Zipping one encoded chunk]
\lean{Complexity.BinArith.zipChunk}
\label{Complexity.BinArith.zipChunk}
\leanok
\uses{Complexity.pairFstD, Complexity.pairSndD, Turing.pairEncode}
Given a function $g$ from Boolean words to Boolean words and a Boolean word $S$, return
the empty word if either the first component of the second component of $S$ or the
second component of the second component of $S$ is empty. Otherwise, pair-encode the
result of applying $g$ to the pair encoding of the first component of $S$ and the pair
encoding of the first components of the first and second components of the second
component of $S$, with the empty word.
\end{definition}

\begin{definition}[Zipping across a word]
\lean{Complexity.BinArith.zipWithW}
\label{Complexity.BinArith.zipWithW}
\leanok
\uses{Complexity.BinArith.zipChunk, Complexity.BinArith.zipStep}
Given a function $g$ from Boolean words to Boolean words and a Boolean word $S$,
concatenate the results of zipping one encoded chunk over the iterates of the zip step
applied to $S$, for iteration indices from $0$ through $|S|+1$.
\end{definition}

\begin{theorem}[Flattening paired decoding rounds]
\lean{Complexity.BinArith.flatMap_zipLoop}
\label{Complexity.BinArith.flatMap_zipLoop}
\leanok
\uses{Complexity.BinArith.decWords, Complexity.BinArith.decWords_of_ne, Complexity.BinArith.length_pairSndD_lt, Complexity.pairFstD, Complexity.pairSndD, Turing.pairDecode}
\difficulty{5}
Let $c$ be any function from pairs of finite bit strings to finite bit strings, and let
$N$ be a natural number. For finite bit strings $u,v$ with $|u|\le N$, consider each
$i=0,\ldots,N-1$: if either string obtained by applying the total second projection $i$
times to $u$ and $v$, respectively, is empty, contribute nothing; otherwise apply $c$ to
the total first projections of those two strings. The concatenation of these
contributions equals the concatenation of $c$ applied to corresponding terms of the
recursively decoded-word lists of $u$ and $v$, stopping when either list runs out.
\end{theorem}

\begin{theorem}[Zip over two encoded word lists]
\lean{Complexity.BinArith.zipWithW_pairEncode}
\label{Complexity.BinArith.zipWithW_pairEncode}
\leanok
\uses{Complexity.BinArith.decWords, Complexity.BinArith.encWords, Complexity.BinArith.flatMap_zipLoop, Complexity.BinArith.iterate_zipStep, Complexity.BinArith.zipChunk, Complexity.BinArith.zipWithW, Complexity.pairFstD_pairEncode, Complexity.pairSndD_pairEncode, Turing.length_pairEncode, Turing.pairEncode}
\difficulty{3}
Let $g$ be any function from finite bit strings to finite bit strings, and let
$q,l_1,l_2$ be finite bit strings. Encode $q$ together with the encodings of $l_1$ and
$l_2$ as a nested pair. Applying $zipWithW$ to this encoding yields the concatenation of
the encodings with empty second component of the outputs $g(\langle q,x,y\rangle)$,
where $(x,y)$ ranges over the termwise pairs of words decoded from $l_1$ and $l_2$,
stopping when either decoded list ends.
\end{theorem}

\begin{theorem}[Polynomial-time repeated pairwise transformation]
\lean{Complexity.BinArith.polyTimeComputable_zipWithW}
\label{Complexity.BinArith.polyTimeComputable_zipWithW}
\leanok
\uses{Complexity.BinArith.iterate_zipStep, Complexity.BinArith.length_iterate_pairSndD_le, Complexity.BinArith.polyTimeComputable_slot2, Complexity.BinArith.polyTimeComputable_slot3, Complexity.BinArith.zipChunk, Complexity.BinArith.zipStep, Complexity.BinArith.zipWithW, Complexity.PolyTimeComputable, Complexity.PolyTimeComputable.comp, Complexity.PolyTimeComputable.pairEncode, Complexity.length_pairSndD_le, Complexity.length_pair_components_le, Complexity.pairFstD, Complexity.pairSndD, Complexity.polyTimeComputable_const, Complexity.polyTimeComputable_emitIter, Complexity.polyTimeComputable_isNil, Complexity.polyTimeComputable_ite, Complexity.polyTimeComputable_or, Complexity.polyTimeComputable_pairFstD, Complexity.polyTimeComputable_pairSndD, Turing.length_pairEncode, Turing.pairEncode}
\difficulty{5}
Let $g$ be a polynomial-time computable function on bit strings. The function that
iterates the specified pairwise transformation $|S|+2$ times on an input $S$, applies
the associated chunk function at each iterate for indices from $0$ through $|S|+1$, and
concatenates the resulting chunks is polynomial-time computable.
\end{theorem}

\begin{definition}[Counting nonempty iterated suffixes]
\lean{Complexity.BinArith.countW}
\label{Complexity.BinArith.countW}
\leanok
\uses{Complexity.pairSndD}
For a Boolean word $w$, form a Boolean word with one true entry for each iteration index
from $0$ through $|w|+1$ at which applying the second-component operation that many
times to $w$ gives a nonempty word; concatenate these entries in index order.
\end{definition}

\begin{theorem}[Unary count of decoded words]
\lean{Complexity.BinArith.countW_eq}
\label{Complexity.BinArith.countW_eq}
\leanok
\uses{Complexity.BinArith.countW, Complexity.BinArith.decWords, Complexity.BinArith.flatMap_listLoop}
\difficulty{4}
For every bit string $w$, the count produced by testing, for each $i$ from $0$ through
$|w|$, whether the second component remains nonempty after $i$ successive
second-component projections is the unary word consisting of one true bit for each word
in the recursive decomposition of $w$. This decomposition is empty for the empty word
and otherwise begins with the first component of $w$, followed by the decomposition of
its second component.
\end{theorem}

\begin{theorem}[Polynomial-time computability of the count function]
\lean{Complexity.BinArith.polyTimeComputable_countW}
\label{Complexity.BinArith.polyTimeComputable_countW}
\leanok
\uses{Complexity.BinArith.countW, Complexity.BinArith.length_iterate_pairSndD_le, Complexity.PolyTimeComputable, Complexity.polyTimeComputable_const, Complexity.polyTimeComputable_emitIter, Complexity.polyTimeComputable_isNil, Complexity.polyTimeComputable_ite, Complexity.polyTimeComputable_pairSndD}
\difficulty{3}
The function that, on a bit string $w$, returns a string consisting of one $1$ for each
iterate $i$ from $0$ through $|w|$ for which the iterated total second projection of $w$
is nonempty is polynomial-time computable.
\end{theorem}
